\subsection{\texorpdfstring{\(\mathbf{R}^{n}\) 与 \(\mathbf{Z}^{n}\) 中的可积性判据}{\textbackslash mathbf\{R\}\^{}\{n\} 与 \textbackslash mathbf\{Z\}\^{}\{n\} 中的可积性判据}}

命题 3. --- 设 N 是 \(\mathbf{R}^{n}\) 上的一个范数，\(r\) 是实数，\(p \in [1, +\infty[\) 。

\begin{enumerate}
\def\labelenumi{\alph{enumi})}
\tightlist
\item
  函数 \((1 + \mathrm{N})^{r}\) 属于 \(\mathcal{L}^{p}(\mathbf{R}^{n}, \mu)\) 当且仅当 \(rp < -n\)；
\end{enumerate}

a') 在 \(\mathbf{Z}^{n}\) 上，函数 \((1+N)^{r}\) 的限制属于 \(\ell^{p}(\mathbf{Z}^{n})\) 当且仅当 \(rp < -n\)；

\begin{enumerate}
\def\labelenumi{\alph{enumi})}
\setcounter{enumi}{1}
\tightlist
\item
  设 V 是 \(\mathbf{R}^{n}\) 中 0 的有界可测邻域。函数 \(\mathbf{N}^{r}\) 属于 \(\mathcal{L}^{p}(\mathbf{R}^{n}-\mathbf{V},\mu)\) 当且仅当 \(rp<-n\) 。
\end{enumerate}

由 \(\mathbf{R}^{n}\) 上范数的等价性（EVT I, p. 14，定理 2 及 TG, IX, p. 32，命题 8），可以假设范数 N 由 \(\mathrm{N}(x)=\sup|x_{i}|\) 给出，其中 \(x=(x_{i})\in\mathbf{R}^{n}\) 。用 B 表示 \(\mathbf{R}^{n}\) 对此范数的单位球。

设 V 是 \(\mathbf{R}^n\) 中 0 的有界可测邻域。函数 \(N^r\) 在 \((\mathrm{B - V})\cup (\mathrm{V - B})\) 上连续且有界，这表明断言 \(b)\) 成立当且仅当 \(V = B\) 时它成立，以下我们就这样假设。函数 \((1 + N)^r\) 在 B 上连续且有界，且满足不等式 \(N^r\leqslant (1 + N)^r\leqslant 2^r N^r\) （在 \(\mathbf{R}^n -\mathbf{B}\) 上），由此可看出断言 \(a)\) 与 \(b)\) 等价。

我们来证明 V = B 时的 b)。可以假设 p = 1。\(r > 0\) 的情形是初等的，故设 \(r \leqslant 0\) 。对每个整数 \(j \geqslant 1\) ，可积集

\[
\mathrm{C} _ {j} = \{x \in \mathbf {R} ^ {n} \mid 2 ^ {j - 1} \leqslant \mathrm{N} (x) <   2 ^ {j} \}
\]

的测度为 \(2^{nj}(2^n - 1)\) 。诸集合 \((\mathrm{C}_j)_{j \geqslant 1}\) 构成 \(\mathbf{R}^n - \mathbf{B}\) 的一个划分。依据 INT, V, p. 4, § 1, no. 1，推论，我们因此有

\[
\begin{array}{c} \int_ {\mathbf {R} ^ {n} - \mathrm{B}} ^ {*} \mathrm{N} ^ {r} d \mu = \sum_ {j \geqslant 1} \int_ {\mathrm{C} _ {j}} ^ {*} \mathrm{N} ^ {r} d \mu \\ \leqslant \sum_ {j \geqslant 1} 2 ^ {n + n j} 2 ^ {r (j - 1)} = 2 ^ {n - r} \sum_ {j \geqslant 1} 2 ^ {(n + r) j} \end{array}
\]

它在 \(r < -n\) 时有限。另一方面，我们有（同前）

\[
\int_ {\mathbf {R} ^ {n} - \mathrm{B}} ^ {*} \mathrm{N} ^ {r} d \mu \geqslant \sum_ {j \geqslant 1} 2 ^ {r j} 2 ^ {n j} (2 ^ {n} - 1) = (2 ^ {n} - 1) \sum_ {j \geqslant 1} 2 ^ {(r + n) j}
\]

它在 \(r \geqslant -n\) 时无穷。

我们类似地证明 \(a'\) )，所考虑的集合 \(\mathrm{C}_{j} \cap \mathbf{Z}^{n}\) 覆盖 \(\mathbf{Z}^{n} - \{0\}\) 并满足

\[
\operatorname{Card} \left(\mathrm{C} _ {j} \cap \mathbf {Z} ^ {n}\right) = (2 ^ {j + 1} - 1) ^ {n} - (2 ^ {j} - 1) ^ {n},
\]

该数属于区间 \([(1-2^{-n})(2^{j+1}-1)^{n},2^{n(j+1)}]\) 。

